\section{The Remaining Days of the Octave: St.~Thomas of Canterbury and St.~Sylvester}\label{sec:octave-days}

After the three feasts of the \textit{Comites Christi}, the calendar opens into
three quieter days before the Octave closes: 29, 30, and 31~December, the
Fifth, Sixth, and Seventh Days within the Octave of the Nativity. Under the
1960 \textit{Codex Rubricarum} (General Rubrics~67--68 and~91, row~17), these
are second-class liturgical days in white vestments. Because St.~Thomas of
Canterbury (29~December) and Pope St.~Sylvester~I (31~December) were reduced to
commemorations in 1960 (\textit{Variationes} n.~5), every non-Sunday weekday
from 29 through 31~December celebrates the same Christmas Octave Mass.

\subsection{The Hybrid Octave Mass and a Present Mystery}

The 1962 Missal crafts a distinctive Mass for 29, 30, and 31~December
(\textit{Missa Puer natus est nobis}, printed pp.~28--29) by joining two of
Christmas Day's own liturgies together. The chants and prayers come from the
Third Mass in the Day---the Introit \textit{Puer natus est}, the Collect
\textit{Concede, quaesumus}, the Gradual from Psalm~97~[98], the Alleluia
\textit{Dies sanctificatus}, the Secret \textit{Oblata}, and the Postcommunion
\textit{Praesta}---while the Epistle and Gospel come from the Second Mass at
Dawn: Titus~3:4--7 (``when the goodness and kindness of God our Saviour
appeared'') and Luke~2:15--20 (the shepherds hurrying to Bethlehem and Mary
pondering all these words in her heart). With the \textit{Gloria}, the Creed,
the Nativity Preface, and the proper \textit{Communicantes} (\textit{diem
sacratissimum}), the Mass invites the Church to linger beside the manger rather
than rush on.

Reading from Pope St.~Leo the Great's Ninth Sermon on the Nativity (Sermon~29,
appointed within the Octave in the Roman Breviary), the Church explains why she
lingers for eight days over a single event:

\begin{quote}
He cometh nearest to the knowledge of the truth, who, the farther he advanceth,
seeth all the more clearly that he can never overtake that whereafter he
searcheth... that we may see the Birth of the Lord, wherein the Word was made
Flesh, not so much as a thing past, as a thing present.
\end{quote}

Because no mind can exhaust the Incarnation in one day, the Octave holds the
Birth of the Lord before our eyes ``not so much as a thing past, as a thing
present.''

\subsection{29 December: St.~Thomas of Canterbury, Bishop and Martyr}

On weekdays from Monday through Saturday, the Fifth Day within the Octave adds
in Low Mass the ordinary commemoration of \textbf{St.~Thomas Becket},
Archbishop of Canterbury and Martyr (omitted in non-conventual Sung Mass, and
omitted altogether when 29~December falls on the Sunday within the Octave).

\begin{massbox}{Collect of St.~Thomas of Canterbury (29 December)}
\textbf{Latin} (1962 Roman Missal). Deus, pro cujus Ecclesia gloriosus Pontifex
Thomas gladiis impiorum occubuit: praesta, quaesumus; ut omnes, qui ejus
implorant auxilium, petitionis suae salutarem consequantur effectum. Per Dominum
nostrum Jesum Christum, Filium tuum: Qui tecum vivit et regnat in unitate.

\textbf{English} (Lasance, 1945). O God, for Whose Church the glorious Bishop
Thomas fell by the swords of wicked men, grant, we beseech Thee, that all who
implore his help may obtain the effect of their petition leading to salvation.
Through our Lord.
\end{massbox}

St.~Thomas entered the Roman calendar in the thirteenth century, carrying the
martyrdom of the \textit{Comites Christi} out of the first century into the
medieval Church's struggle for the freedom of the Church (\textit{libertas
Ecclesiae}). After years of exile in France defending ecclesiastical liberties
against King Henry~II, Thomas returned to Canterbury in December~1170 knowing
death stood near. In the universal 1962 Breviary, 29~December keeps the Office
of the Fifth Day within the Octave with only a Lauds commemoration of
St.~Thomas; yet the traditional pre-1960 Roman Breviary preserves in its Sixth
Lesson the vivid account of the evening of 29~December~1170 inside Canterbury
Cathedral, when armed knights burst in during Vespers:

\begin{quote}
The clergy in attendance on him, conscious of the attempt about to be made,
wished to bolt the doors. But the saint caused them to be again opened, saying,
``The Church of God is not to be made a castle of, and for the cause of God's
Church I am willing to die.'' He then said to his murderers, ``I charge you in
the name of the Almighty God to hurt none of my people.''
\end{quote}

Struck down beside the altar, he was canonized by Pope Alexander~III in~1173.
In the Third Nocturn of that older Office, St.~John Chrysostom's Homily~59 on
John (John~10:11--16) explains the difference between the Good Shepherd and the
hireling: when the wolf comes, the hireling flees because the sheep are not his
own; the true pastor stays and lays down his life. Standing four days after
Bethlehem, St.~Thomas reminds the Church that the Child laid in a manger is the
Lord of a free kingdom no earthly crown can buy or intimidate.

\subsection{31 December: Pope St.~Sylvester~I and the Turning of the Year}

On 30~December (the Sixth Day within the Octave), no saint is commemorated;
the day belongs wholly to the quiet of the Octave (or to the Sunday within the
Octave when 25~December falls on Tuesday). On 31~December (the Seventh Day
within the Octave), Low Mass on Monday through Saturday adds the ordinary
commemoration of \textbf{Pope St.~Sylvester~I}, Confessor (omitted in
non-conventual Sung Mass and whenever 31~December is the Sunday within the
Octave), whose ancient Roman station stood at his burial place in the Cemetery
of Priscilla on the Via Salaria.

\begin{massbox}{Collect of St.~Sylvester~I (31 December)}
\textbf{Latin} (1962 Roman Missal). Gregem tuum, Pastor aeterne, placatus
intende: et per beatum Silvestrum Summum Pontificem perpetua protectione
custodi; quem totius Ecclesiae praestitisti esse pastorem. Per Dominum nostrum.

\textbf{English} (Lasance, 1945, Common of Supreme Pontiffs). Eternal Shepherd,
do Thou look favorably upon Thy flock, which we beseech Thee to guard and keep
for evermore through the Blessed [Sylvester], Supreme Pontiff, whom Thou didst
choose to be the chief shepherd of the whole Church. Through our Lord.
\end{massbox}

Pope St.~Sylvester governed the Church of Rome from 314 to 335, across the
epochal passage from imperial persecution to the Constantinian peace. During
his pontificate the great Roman basilicas---the Lateran, St.~Peter's,
St.~Paul's, St.~Lawrence, and Holy Cross in Jerusalem---rose above the martyrs'
graves, and the \textbf{First Council of Nicaea (325)} confessed the eternal
Son to be consubstantial (\textit{homoousios}) with the Father against the
Arian heresy. While later Roman legend pictured Sylvester baptizing Emperor
Constantine in Rome (introduced cautiously in the older pre-1960 Breviary as ``an
old story in the Church of Rome,'' whereas Constantine was in fact baptized on
his deathbed at Nicomedia in~337), Sylvester's real achievement stands at the
very heart of Christmas: the Roman see that kept the earliest documented feast
of 25~December (Section~\ref{sec:history}) and defended the full divinity of
the Child of Bethlehem at Nicaea.

In the universal 1962 Breviary, 31~December likewise keeps the Office of the
Seventh Day within the Octave with only a Lauds commemoration of St.~Sylvester;
yet the Sixth Lesson of his traditional pre-1960 Office preserves a pastoral
memory ideally suited to New Year's Eve: St.~Sylvester confirmed the Church's
liturgical custom of naming the weekdays \textbf{\textit{Feria}} rather than
calling them after pagan planetary gods, a word that ``signifieth an `holiday,'
and pointeth to the duty of the clergy ever to lay aside all worldly labour,
and leave themselves free to do continually the work of the Lord.''

At dusk on 31~December, as First Vespers of the Octave Day opens and the civil
year draws to its close, Catholic churches and households have long sung or
recited the \textit{Te Deum} in thanksgiving for every mercy of the departing
year (as Bl.~Ildefonso Schuster notes in his commentary on St.~Sylvester's
day). Paired with the fifth part of St.~Francis de Sales's \textit{Introduction
to the Devout Life}---an unhurried annual examination of conscience and renewal
of good resolutions before God---the evening of 31~December turns the hinge of
the calendar from noisy distraction into grateful adoration.
